\section{Cooling Tower}

\subsection{Mathematical Problem Formulation and Geometric Parametrisation}

The cooling tower of a nuclear power plant defines a continuous surface optimisation problem over a rotationally symmetric structure. 
To render this tractable computationally, the continuous domain is discretised into a vertical stack of $m$ conical frustums, bounded by $m+1$ horizontal cross-sections (rings). 
Let $z$ denote the vertical axis originating at the tower base, with $z_0 = 0$. 
Each frustum $i \in \{1, \dots, m\}$ is fully characterised by its bottom radius $r_{i-1}$, top radius $r_i$, and section height $h_i$.

The macroscopic engineering objective is to minimise construction material costs, which are directly proportional to the total lateral surface area $A_{\text{total}}$. 
This is subject to a strict equality constraint enforcing a fixed thermodynamic cooling capacity, expressed as a target volume $V_{\text{target}} = 70{,}320\,\text{m}^3$. 
For each frustum $i$, the slant height $s_i$, lateral surface area $A_i$, and volumetric capacity $V_i$ are given by:

\begin{align}
    s_i &= \sqrt{(r_{i-1} - r_i)^2 + h_i^2} \\
    A_i &= \pi (r_{i-1} + r_i)\, s_i \\
    V_i &= \frac{\pi h_i}{3} \left( r_{i-1}^2 + r_{i-1} r_i + r_i^2 \right)
\end{align}

The resulting global problem is formally a constrained nonlinear programme (NLP):

\begin{equation}
    \min_{\mathbf{x} \in \mathcal{D}}\; A_{\text{total}}(\mathbf{x}) = \sum_{i=1}^{m} A_i(\mathbf{x})
    \quad \text{subject to} \quad
    \mathcal{H}(\mathbf{x}) = \sum_{i=1}^{m} V_i(\mathbf{x}) - V_{\text{target}} = 0
\end{equation}

where $\mathbf{x}$ is the decision variable vector, and $\mathcal{D} \subseteq \mathbb{R}^{D}$ is the feasible space defined by the physical boundary conditions.

\subsection{Methodology: Quadratic Penalty Formulation and Search Landscape}

To evaluate the three fundamentally different algorithms, Simulated Annealing (SA), Particle Swarm Optimisation (PSO), and the quasi-Newton BFGS method, on a common, unbiased basis, the equality constraint $\mathcal{H}(\mathbf{x}) = 0$ must be handled uniformly. 
Since SA and PSO are natively unconstrained, and since introducing solver-specific constraint-handling logic would compromise benchmark fairness, the equality constraint is absorbed into the objective via a \textbf{quadratic penalty method}. 

The surrogate objective becomes:

\begin{equation}
    J(\mathbf{x}) = A_{\text{total}}(\mathbf{x}) + \rho \left[\mathcal{H}(\mathbf{x})\right]^2
\end{equation}

where $\rho \gg 0$ is a static scalar penalty weight. A static $\rho$ is deliberately chosen over an adaptive schedule to ensure that all three algorithms operate on an identical, fixed search landscape any dynamic $\rho$ strategy would introduce algorithm-dependent trajectory effects that would confound the performance comparison. 
The chosen value $\rho = 10^3$ was calibrated empirically to ensure that volume violations incur a penalty several orders of magnitude larger than the surface area objective, effectively rendering infeasible solutions non-competitive without distorting the local geometry near the feasible manifold.

This penalty formulation has two notable consequences on the search landscape. 
First, for BFGS, the quadratic term degrades the conditioning of the Hessian approximation in proportion to $\rho$, increasing the curvature of the surrogate near constraint boundaries~\cite{nocedal2006numerical}. 
Second, for SA and PSO, the penalty creates steep regions of high curvature in the search landscape around infeasible points, placing greater demands on the exploration-exploitation balance of each algorithm.

\subsection{Architectural Decoupling and Benchmark Fairness}

A fully decoupled, object-oriented software architecture was implemented around a central \texttt{CoolingTowerProblem} class, which encapsulates the geometric model and all case-specific logic. 
The key design principle is a vector mapping pipeline (\texttt{unpack\_decision\_variables}): all three solvers interact exclusively with an abstract decision vector $\mathbf{x} \in \mathbb{R}^D$, with no knowledge of the underlying physical structure. 
The class internally reconstructs the full radii and heights arrays from $\mathbf{x}$ according to the active case. 
This architecture ensures that the \texttt{solver.py} orchestrator applies each algorithm identically across all eight cases, with no solver-specific preprocessing or constraint logic injected at the call level. 
The result is a genuinely fair, one-to-one algorithmic comparison.

\subsection{Study Cases: Dimensionality and Physical Complexity}

Eight cases were designed to progressively stress-test the algorithms across increasing levels of search space complexity, physical realism, and parametric dimensionality. 
All cases are evaluated against BFGS, SA, and PSO.

\begin{itemize}

    \item \textbf{Case 1: Mandatory Reference Case ($D = 9$):}
    The canonical benchmark. Section heights $h_i$ and boundary radii $r_0 = 39.3\,\text{m}$, $r_m = 27.4\,\text{m}$ are fixed, and the nine inner radii $r_1, \dots, r_9$ constitute the sole decision variables. 
	This low-dimensional, well-conditioned problem serves as the convergence baseline against which all other cases are evaluated.

    \item \textbf{Case 2: Variable Heights ($D = 10$):}
    The radii profile from Case 1 is fixed, and the ten section heights $h_i$ are optimised instead. 
	This isolates the vertical sensitivity of the volume constraint and characterises how each algorithm navigates a height-dominated feasible manifold.

    \item \textbf{Case 3: Full Freedom ($D = 19$):}
    Both inner radii and section heights are optimised simultaneously.
	The near-doubling of dimensionality substantially expands the search landscape, exposing each algorithm to a higher density of local minima and testing resilience against premature convergence.

    \item \textbf{Case 4: Tight Waist ($D = 9$):}
    A structural constraint is imposed: $r_{\text{mid}} < 0.7\,\min(r_0, r_m)$, artificially narrowing the tower throat. 
	This breaks the convexity of the feasible region and introduces a disconnected, non-convex feasible zone, testing each solver's ability to navigate segmented constraint boundaries.

    \item \textbf{Case 5: Cylindrical Initialisation ($D = 9$):}
    A deliberately poor starting point is imposed: $r_i = r_0$ for all $i$, corresponding to a cylinder that substantially violates the volume constraint. 
	This case measures the sensitivity of each algorithm to initialisation, and in particular the speed at which each method escapes a far-from-feasible equilibrium.

    \item \textbf{Case 6: High Volume ($D = 9$):}
    The target volume is increased by $50\%$ to $V_{\text{target}} = 105{,}480\,\text{m}^3$. 
	This rescales the feasible manifold and alters the relative magnitude of the penalty gradient $\nabla[\mathcal{H}(\mathbf{x})]^2$, probing each algorithm's scalability under a shifted constraint landscape.

    \item \textbf{Case 7: Restricted Bounds ($D = 9$):}
    Box constraints are imposed on all inner radii: $r_i \in [20, 60]\,\text{m}$, simulating civil engineering construction tolerances. 
	This tests BFGS boundary approximations via L-BFGS-B, and requires SA and PSO to implement explicit boundary clipping or reflection strategies.

    \item \textbf{Case 8: Hyperbolic Parametrisation ($D = 3$):}
    Rather than optimising discrete node positions, the tower profile is constrained to a hyperboloid of one sheet:
    \begin{equation}
        r(z) = a\sqrt{1 + \frac{(z - c)^2}{b^2}}
    \end{equation}
    where $a$ is the waist radius, $b$ controls the curvature, and $c$ is the height of the waist. The optimisation operates directly on $(a, b, c)$, reducing the search to $D = 3$. This parametrisation substantially mitigates the curse of dimensionality while inherently guaranteeing $C^1$ profile continuity---a structurally significant property, as smooth surface curvature improves wind load distribution along the tower shell~\cite{mungan1976buckling}. This case also enables a direct comparison between the discrete frustum approximation and the analytically integrated surface area of the true hyperboloid, quantifying the geometric discretisation error introduced in Cases~1--7.

\end{itemize}


\section{Algorithms}

\subsection{Particle Swarm Optimisation method}

TO DO

\subsection{Simulated Annealing method}

TO DO

\subsection{Broyden-Fletcher-Goldfarb-Shanno method}

TO DO